% =============================================================================
%  Sección 1: Marco Conceptual de los Humedales Urbanos
% =============================================================================
% Archivo: Marco_Conceptual.tex
% Materia: Desarrollo Urbano Sostenible
% Capítulo 1: Marco Conceptual de los Humedales

\section{Marco Conceptual de los Humedales Urbanos}

En esta sección se explorará el marco normativo y las definiciones técnicas sobre los humedales dentro de la Ciudad de México y el Área Metropolitana
con especial énfasis en los ubicados en la alcaldía de Xochimilco (Humedales de Xochimilco) y Aragón en la Alcaldía de Gustavo A. Madero que colinda
con los municipios de Ecatepec y Nezahualcóyotl en el Estado de México. A su vez, se hará la distinción entre ambos ecosistemas; siendo uno creado de
manera natural (como lo son los humedales de Xochimilco y los ejidos de San Gregorio Atlapulco) y los creados artificialmente en la zona de Aragón que
se encuentran en el Bosque de San Juan de Aragón.

\subsection{¿Qué es un humedal?}
De acuerdo con la Convención de Ramsar de 1971, los humedales se definen técnicamente como ecosistemas diversos que actúan como una zona de transición entre
los sistemas acuáticos y terrestres, abarcando ambientes tales como turberas, pantanos, llanuras de inundación, lagos y ríos \cite{vergara2024humedales}. 
En estos entornos, el agua constituye el factor principal que controla el medio ambiente y la vida vegetal y animal asociada, incluyendo zonas costeras que no 
excedan los seis metros de profundidad en marea baja \cite{paot2006xochimilco, vergara2024humedales}.

Ecológicamente, estos sistemas presentan variaciones en sus cuerpos de agua según la dinámica hidrológica local y climática, caracterizándose por poseer suelos
hídricos y vegetación específicamente adaptada a condiciones de saturación prolongada \cite{vergara2024humedales}. Su importancia en el ámbito urbano radica 
en su alta productividad biológica y en servicios ecosistémicos críticos como la mitigación de inundaciones, la retención de nutrientes y el almacenamiento de 
carbono \cite{vergara2024humedales}.

Para efectos de este estudio, se establece una diferenciación clave basada en el origen y la gestión de estos sistemas en la capital:

\begin{itemize}
    \item \textbf{Humedales Naturales (Remanentes lacustres):} Representan los vestigios del antiguo sistema de lagos de la Cuenca de México.
    El caso de \textbf{Xochimilco} destaca como un humedal permanente constituido por una red de canales y lagunas, reconocido internacionalmente
    desde 2004 como el \textit{Sitio Ramsar 1363} \cite{paot2006xochimilco}. No obstante, este sistema enfrenta un deterioro ambiental constante 
    derivado de la presión de la mancha urbana y la ocupación irregular \cite{paot2006xochimilco, unam2020deterioro}.
    
    \item \textbf{Humedales Artificiales y Construidos (Ingeniería Ecológica):} Son sistemas diseñados para emular procesos naturales con fines 
    de remediación. Un ejemplo emblemático es el humedal del \textbf{Bosque de San Juan de Aragón}, donde se ha implementado tecnología universitaria 
    para mejorar la calidad del agua mediante procesos biológicos. Asimismo, dentro de Xochimilco existen infraestructuras artificiales de control hídrico,
    como las lagunas de regulación Ciénega Chica y Ciénega Grande, destinadas primordialmente a la prevención de inundaciones \cite{paot2006xochimilco}.
\end{itemize}

\subsection{Importancia tridimensional en el ámbito urbano}
Los humedales no solo fungen como un ecosistema importante en el ámbito medio-ambiental. Su relevancia debe ser analizada desde un punto de vista integral 
para su correcto análisis como elementos estratégicos para la resiliencia urbana a través de tres dimensiones fundamentales:
%
\subsubsection*{Dimensión Ambiental}
Los humedales urbanos actúan como reguladores críticos del entorno biofísico. En términos climáticos, desempeñan un papel fundamental en la regulación de
la temperatura y las condiciones climáticas locales \cite{paot2006xochimilco}. Su capacidad para el almacenamiento de carbono es vital para la mitigación del cambio climático,
funcionando como sumideros naturales de gases de efecto invernadero \cite{paot2006xochimilco}.

Desde una perspectiva hidrológica, estos ecosistemas estabilizan las capas freáticas al retener el agua de lluvia y liberarla lentamente, funcionando como reservorios naturales 
que mitigan inundaciones y brindan protección contra tormentas. En el caso del humedal artificial de Aragón, la implementación de ingeniería ecológica ha permitido la depuración 
de aguas tratadas, generando servicios ambientales adicionales para la metrópoli. Asimismo, representan refugios críticos de biodiversidad; mientras Xochimilco alberga especies 
endémicas como el ajolote (\textit{Ambystoma mexicanum}), el humedal de Aragón ha logrado triplicar su avifauna en 20 años, pasando de 60 especies en 2001 a más de 182 en 2022.

%
\subsubsection*{Dimensión Social}
Socialmente, estos espacios ofrecen beneficios directos a la salud pública y el bienestar ciudadano. Funcionan como áreas destinadas a la recreación, la investigación y la educación ambiental, 
destacando el Parque Ecológico de Xochimilco y el Bosque de San Juan de Aragón como nodos de esparcimiento masivo. La presencia de vegetación hidrófila y jardines para polinizadores contribuye a
la mejora de la calidad del aire y la escala de especies migratorias, elevando la calidad ambiental urbana.

En el ámbito histórico, el sistema de \textbf{chinampas} en Xochimilco representa un paisaje cultural único en el mundo, reconocido por la UNESCO en 1987 como Patrimonio Mundial Cultural y Natural.
Este sistema agrícola prehispánico no solo preserva la identidad histórica de la Cuenca de México, sino que constituye un legado sociocultural que articula la vida comunal y productiva de los pobladores locales.

\subsubsection*{Dimensión Política}
La dimensión política se define por la tensión entre la protección jurídica y la presión del desarrollo urbano. El marco de gobernanza está delimitado por declaratorias internacionales y locales, como el 
\textbf{Sitio Ramsar 1363} y el Área Natural Protegida (ANP) decretada en 1992. Estos instrumentos obligan al Estado a asegurar el mantenimiento de las características ecológicas frente a la expansión de la 
infraestructura gris.

El reto principal radica en la gestión de políticas públicas frente a la presión inmobiliaria y la construcción de infraestructura vial. Un caso emblemático es la orden judicial de paralizar el puente vehicular 
en Periférico Sur y Canal Nacional (Cuemanco), donde un juez federal determinó que la protección del Sitio Ramsar prevalece sobre el proyecto urbano, con el fin de evitar daños irreparables al humedal. 
Este escenario político resalta la necesidad de una gobernanza que priorice soluciones basadas en la naturaleza frente al deterioro ambiental derivado de la ocupación irregular y la fragmentación del territorio.
